\documentclass[11pt,a4paper]{article}
\usepackage[utf8]{inputenc}
\usepackage[vietnamese]{babel}
\usepackage{amsmath,amssymb}
\usepackage{geometry}
\geometry{a4paper, total={180mm,260mm}, left=15mm, right=15mm, top=18mm, bottom=18mm}
\usepackage{tikz}
\usetikzlibrary{shapes.geometric, arrows.meta, positioning, calc}
\usepackage{tcolorbox}
\tcbuselibrary{skins,breakable}
\usepackage{xcolor}
\usepackage{multicol}
\usepackage{array}
\usepackage{fontawesome}
\usepackage{fancyhdr}

\pagestyle{fancy}
\fancyhf{}
\renewcommand{\headrulewidth}{0.5pt}
\renewcommand{\footrulewidth}{0.5pt}
\fancyhead[L]{\small\bfseries\color{teal}\faBook\ SỔ TAY KIẾN THỨC TOÁN 3}
\fancyhead[R]{\small\bfseries\color{teal}Lớp Toán Thầy Lĩnh -- loptoanthaylinh.vn}
\fancyfoot[C]{\small\color{gray}-- Sổ tay kiến thức: Bài 31. Gam (g) --}

\definecolor{primary}{HTML}{005F73}
\definecolor{secondary}{HTML}{0A9396}
\definecolor{accent}{HTML}{EE9B00}
\definecolor{danger}{HTML}{AE2012}
\definecolor{boxbg}{HTML}{F4F6F8}

\begin{document}

% --- BANNER TIÊU ĐỀ ---
\begin{tcolorbox}[
  enhanced,
  colback=white,
  colframe=primary,
  arc=3mm,
  boxrule=1.2pt,
  left=4mm, right=4mm, top=2.5mm, bottom=2.5mm,
  drop shadow=black!10!white
]
  \centering
  {\Large\bfseries\color{primary}\faBookmark\ BÀI 31: GAM (g)}\\[1.5mm]
  {\small\itshape\color{black!75}Đơn vị đo khối lượng, quan hệ 1 kg = 1000 g và kĩ năng đọc cân đồng hồ, cân đĩa trên 1 trang A4}
\end{tcolorbox}

\vspace{-2mm}


\noindent
\begin{minipage}[t]{0.49\textwidth}\vspace{0pt}
\begin{tcolorbox}[
  enhanced,
  colback=blue!4!white,
  colframe=blue!75!black,
  arc=2mm, boxrule=0.8pt,
  left=2mm, right=2mm, top=1.8mm, bottom=1.8mm,
  title={\footnotesize\bfseries\color{white}\faKey\ 1. Đơn vị Gam (g)}
]
{\footnotesize
\begin{itemize}\setlength{\itemsep}{1.5pt}\setlength{\leftskip}{-2.5mm}
  \item Gam là đơn vị đo khối lượng, viết tắt là \textbf{g}.
  \item Thường dùng để cân các vật nhẹ (quả cam, gói bánh, bao thư, đồ trang sức...).
  \item \textbf{Mối quan hệ vàng với Ki-lô-gam:}
  $$\mathbf{1\text{ kg} = 1000\text{ g}} \iff \mathbf{1000\text{ g} = 1\text{ kg}}$$
  \item \textbf{Các phân số khối lượng quen thuộc:}
  \begin{itemize}\setlength{\itemsep}{0.5pt}\setlength{\leftskip}{-2.5mm}
    \item Nửa ki-lô-gam: $\dfrac{1}{2}\text{ kg} = 1000 : 2 = \mathbf{500\text{ g}}$.
    \item Một phần tư ki-lô-gam: $\dfrac{1}{4}\text{ kg} = \mathbf{250\text{ g}}$.
  \end{itemize}
\end{itemize}
}
\end{tcolorbox}
\end{minipage}
\hfill
\begin{minipage}[t]{0.49\textwidth}\vspace{0pt}
\begin{tcolorbox}[
  enhanced,
  colback=teal!4!white,
  colframe=teal!75!black,
  arc=2mm, boxrule=0.8pt,
  left=2mm, right=2mm, top=1.8mm, bottom=1.8mm,
  title={\footnotesize\bfseries\color{white}\faBalanceScale\ 2. Kĩ năng đọc cân các loại}
]
{\footnotesize
\begin{itemize}\setlength{\itemsep}{1.5pt}\setlength{\leftskip}{-2.5mm}
  \item \textbf{Cân đồng hồ:} Kim chỉ vào vạch nào thì vật nặng bằng số ghi ở vạch đó (chú ý đơn vị trên mặt cân là $\text{kg}$ hay $\text{g}$).
  \item \textbf{Cân hai đĩa (Cân Rô-béc-van):}
  Khi cân thăng bằng:
  $$\text{Khối lượng vật} = \text{Tổng khối lượng quả cân}$$
  \item Nếu quả cân ở cả hai đĩa: Khối lượng vật bằng hiệu khối lượng các quả cân giữa 2 đĩa.
\end{itemize}
}
\end{tcolorbox}
\end{minipage}

\vspace{-2.5mm}
% --- KHỐI 3: SƠ ĐỒ TƯ DUY TIKZ ---
\noindent
\begin{tcolorbox}[
  enhanced,
  colback=orange!4!white,
  colframe=orange!80!black,
  arc=2mm, boxrule=0.8pt,
  left=2mm, right=2mm, top=1.5mm, bottom=1.5mm,
  title={\footnotesize\bfseries\color{white}\faSitemap\ SƠ ĐỒ TƯ DUY: HỆ KHỐI LƯỢNG VÀ KĨ THUẬT ĐỌC CÂN}
]
\centering
\begin{tikzpicture}[
  box/.style={draw, rounded corners=2pt, align=center, font=\scriptsize\bfseries, inner sep=2.5pt, minimum height=0.75cm},
  arr/.style={-{Stealth[length=2mm]}, thick}
]
  \node[box, fill=purple!15, draw=purple!80!black, text width=2.6cm] (g1) at (0, 0) {ĐƠN VỊ GAM\\Viết tắt: $\mathbf{g}$\\Vật khối lượng nhỏ};
  \node[box, fill=teal!15, draw=teal!80!black, text width=2.7cm] (g2) at (3.5, 0) {QUY ĐỔI\\$\mathbf{1\text{ kg} = 1000\text{ g}}$\\$500\text{ g} = \frac{1}{2}\text{ kg}$};
  \node[box, fill=blue!12, draw=blue!80!black, text width=2.6cm] (g3) at (7.0, 0) {CÂN HAI ĐĨA\\Thăng bằng:\\Vật = Tổng quả cân};
  \node[box, fill=orange!15, draw=orange!80!black, text width=2.6cm] (g4) at (10.5, 0) {PHÉP TÍNH\\Cùng đơn vị $\text{g}$\\Ghi kèm đơn vị $\text{g}$};

  \draw[arr, purple!80!black] (g1) -- (g2);
  \draw[arr, teal!80!black] (g2) -- (g3);
  \draw[arr, blue!80!black] (g3) -- (g4);
\end{tikzpicture}
\end{tcolorbox}

\vspace{-2.5mm}
% --- KHỐI 4: VÍ DỤ & MẸO TRÁNH BẪY ---
\noindent
\begin{minipage}[t]{0.49\textwidth}\vspace{0pt}
\begin{tcolorbox}[
  enhanced,
  colback=green!4!white,
  colframe=green!60!black,
  arc=2mm, boxrule=0.8pt,
  left=2mm, right=2mm, top=1.8mm, bottom=1.8mm,
  title={\footnotesize\bfseries\color{white}\faCheckCircle\ Ví dụ thực hành chuẩn mực}
]
{\footnotesize
\textbf{1. Đổi đơn vị khối lượng:}\\
$2\text{ kg} = \mathbf{2000\text{ g}}$.\\
$1\text{ kg } 250\text{ g} = 1000\text{ g} + 250\text{ g} = \mathbf{1250\text{ g}}$.\\
$3000\text{ g} = \mathbf{3\text{ kg}}$.\\
\textbf{2. Bài toán cân đĩa:}\\
Đĩa trái đặt quả dưa hấu, đĩa phải đặt quả cân $2\text{ kg}$ và $500\text{ g}$. Cân thăng bằng. Quả dưa hấu nặng:\\
$2\text{ kg} + 500\text{ g} = \mathbf{2500\text{ g}}$ (hoặc $2\text{ kg } 500\text{ g}$).
}
\end{tcolorbox}
\end{minipage}
\hfill
\begin{minipage}[t]{0.49\textwidth}\vspace{0pt}
\begin{tcolorbox}[
  enhanced,
  colback=red!4!white,
  colframe=red!70!black,
  arc=2mm, boxrule=0.8pt,
  left=2mm, right=2mm, top=1.8mm, bottom=1.8mm,
  title={\footnotesize\bfseries\color{white}\faExclamationTriangle\ Mẹo tránh bẫy trắc nghiệm}
]
{\footnotesize
\begin{itemize}\setlength{\itemsep}{1.5pt}\setlength{\leftskip}{-2.5mm}
  \item \textbf{Bẫy cộng không cùng đơn vị:}\\
  $1\text{ kg} + 200\text{ g} \ne 201\text{ kg}$ hay $201\text{ g}$! Phải đổi: $1000\text{ g} + 200\text{ g} = \mathbf{1200\text{ g}}$.
  \item \textbf{Bẫy cân hai đĩa có quả cân hai bên:}\\
  Đĩa trái có vật và quả cân $200\text{ g}$, đĩa phải có quả cân $1\text{ kg}$.\\
  Khối lượng vật là: $1000 - 200 = \mathbf{800\text{ g}}$, không được cộng vào thành $1200\text{ g}$!
\end{itemize}
}
\end{tcolorbox}
\end{minipage}


\end{document}
